{\large\textbf{QUESTION T3 [10 pts]. Chasing the absolute zero}}

A refrigeration cycle must be carried out by a system with characteristics
appropriate for the temperatures to be achieved. One hundred years ago, in
1926, the chemist William Giauque and, independently, the physicist Peter Debye
proposed paramagnetic salts as ideal systems for lowering temperatures far
below $1\,\mathrm{K}$. Indeed, beginning in the following decade, temperatures
as low as $0.0044\,\mathrm{K}$ were achieved (Wander de Haas and Gerrit
Wiersma, 1935, Leiden Cryogenic Laboratories). In this question, the
paramagnetic Carnot refrigeration cycle is studied.

\begin{center}\includegraphics[width=0.40\linewidth]{figures/IPhO_2026_Q3_fig1.png}\end{center}

\begin{center}\texttt{Fig. 3a.}\end{center}

Consider a torus with mean radius $R$ and inner radius $r$ made of a
homogeneous, isotropic paramagnetic material, around which an insulated
conducting wire is wound, with its ends connected to an external voltage source
(emf), as shown in figure 3a. The resistance of the wire is so low that energy
losses due to the heating of the wire can be neglected. The winding is dense,
with $N$ turns in total, and it is assumed that $r \ll R$, so that the fields
$\vec{H}$ and $\vec{B}$ and the magnetization $\vec{M}$ have approximately
constant magnitudes throughout the torus. Let $V$ and $A$ be the volume and the
cross-section area of the torus. Hereinafter, the abbreviation Pm-T is used to
refer to the paramagnetic torus.

In paramagnetic materials, $\vec{M}$ is parallel to $\vec{H}$ and the following
relation holds:
\[
\overrightarrow{B} = \mu_0\overrightarrow{H} + \mu_0\overrightarrow{M}.
\]

\textbf{Energy transfer sign convention:} in this problem, Work ($W$) and Heat
($Q$) are considered positive when they flow into the system, and negative when
they flow out.

\textbf{T3-A [1.0 pts]. Work in the Pm-T}

\textit{Hint: in magnetic materials, $\oint_C \vec{H}\cdot d\vec{\ell} = I_C$,
where $I_C$ is the net free current passing through the area bounded by the
closed curve $C$.}

\textbf{T3-A1 [0.2 pts].} Let $H$ be the magnitude of $\vec{H}$ in the torus.
Write $H$ in terms of $N$, $A$, $V$ and the instantaneous electric current $I$
in the wire.

\textbf{T3-A2 [0.6 pts].} Let $dW_{emf}$ be the work that the external voltage
source performs to change the magnitude of $\vec{B}$ by $dB$. Write $dW_{emf}$
in terms of $V$, $H$, and $dB$.

\textbf{T3-A3 [0.2 pts]}. The total work done by the voltage source,
$dW_{emf}$, can be divided in two parts: the work that would be performed to
change the magnetic field if the toroid had a vacuum core, $dW_{vac}$, and the
work done on the paramagnetic material itself, $dW$. Write $dW$ in terms of
$\mu_0$, $H$, $V$ and $dM$.

\textbf{T3-B [3.0 pts]. Heat and temperature in the Pm-T}

The following table compares two thermodynamic systems: the ideal gas and the
Pm-T. Use this information and your answer to question T3-A3 to solve the
following questions. This information will also be used in part C. If you did
not find an answer for $dW$ in question T3-A3, you may use
$dW = \alpha V\mu_0 HdM$ where $\alpha$ is a number. The volume $V$ of the
paramagnetic torus is constant, so that there is no work due to pressure.
Below, $T$ is the temperature in all instances.

\begin{center}
\begin{tabular}{|c|c|}
\hline
\textbf{IDEAL GAS} & \textbf{PARAMAGNETIC TORUS} \\
\hline
State functions: $P, V$ and temperature $T$ & State functions: $H, M$ and temperature $T$ \\
\hline
Constants: $R, \gamma$ and number of moles $n$ & Constants: $K, \lambda$ and number of moles $n$ \\
\hline
Equation of state: $PV = nRT$ & Equation of state: $TMV = nKH$ \\
\hline
\begin{tabular}{@{}c@{}}Heat capacity at constant $V$:\\ $C_V = nR/(\gamma - 1)$\end{tabular}
& Heat capacity at constant $M$: $C_M = n\lambda/T^2$ \\
\hline
\begin{tabular}{@{}c@{}}Internal energy $U$ in any process satisfies\\ $dU = C_VdT$\end{tabular}
& \begin{tabular}{@{}c@{}}Internal energy $U$ in any process satisfies\\ $dU = C_MdT$\end{tabular} \\
\hline
\end{tabular}
\end{center}

$K$ and $\lambda$ are constants given by the material in the torus. You may use
the equations provided in the table directly without proving them.

\textbf{T3-B1 [1.5 pts].} If the magnitude of $\vec{H}$ changes from $H_i$ to
$H_f$ at a constant temperature $T$, an energy $Q$ is transferred between the
Pm-T and its environment. Write $Q$ in terms of $\mu_0$, $n$, $V$, $K$, $H_i$,
$H_f$, and $T$.

\textbf{T3-B2 [1.5 pts].} If the magnitude of $\vec{H}$ changes adiabatically
from $H_i$ to $H_f$, the temperature of the Pm-T changes from $T_i$ to $T_f$.
Write $\Delta T = T_f - T_i$ in terms of $\mu_0$, $V$, $n$, $K$, $\lambda$,
$H_i$, $H_f$, and $T_i$.

\textbf{T3-C [6.0 pts]. Carnot refrigerator of the Pm-T}

\begin{center}\includegraphics[width=0.39\linewidth]{figures/IPhO_2026_Q3_fig2.png}\end{center}

\begin{center}Fig.3b.\end{center}

The $H$ vs. $T$ diagram of the Carnot refrigeration cycle
$1 \rightarrow 2 \rightarrow 3 \rightarrow 4 \rightarrow 1$ performed by the
Pm-T is shown in figure 3b. In this cycle, let:

$T_h$: temperature of the hot reservoir\\
$T_c$: temperature of the cold reservoir\\
$Q_h$: absolute value of the heat transferred to the hot reservoir\\
$Q_c$: absolute value of the heat absorbed from the cold reservoir

\textbf{T3-C1 [0.2 pts].} Mark $T_h$ and $T_c$ on the $T$ axis of the $H$ vs.
$T$ diagram. Also label the processes in which the transfers of $Q_h$ and $Q_c$
occur, respectively.

\textbf{T3-C2 [1.5 pts].} Let $M_1$, $M_2$, $M_3$, and $M_4$ be the magnitudes
of $\vec{M}$ at the vertices of this cycle. Write $M_1$ in terms of $M_2$,
$M_3$, and $M_4$.

\textbf{T3-C3 [0.8 pts].} Suppose that to cool $1.00\,\mathrm{L}$ of liquid
helium initially at $1.00\,\mathrm{K}$, we utilize a 2.0 moles Pm-T made of
\textit{potassium chromate}. For this material,
$K = 1.87 \times 10^{-6}\,\mathrm{K{\cdot}m^3/mol}$, its density is
$2730\,\mathrm{kg/m^3}$ and its molar mass is $0.19\,\mathrm{kg/mol}$. The
Pm-T performs a Carnot refrigeration cycle with $H_1 = 411624\,\mathrm{A/m}$,
$H_2 = 311306\,\mathrm{A/m}$, $H_3 = 204618\,\mathrm{A/m}$ and
$H_4 = 240446\,\mathrm{A/m}$. After one operating cycle, calculate the final
temperature of the liquid helium. Assume that for the liquid helium both the
specific heat capacity $c = 100\,\mathrm{J/(kg{\cdot}K)}$ and the density
$\rho = 130\,\mathrm{kg/m^3}$ remain constant.

\textbf{T3-C4 [2.0 pts].} In order to continuously lower the temperature of a
body with constant heat capacity $C_c$, Carnot cycles are allowed to occur over
time in such a way that the power $P$ transferred to the refrigerator remains
constant. It is then assumed that the temperature $T_h$ of the hot reservoir
remains constant and that in each operating cycle the temperature $T_c$ of the
cooled body decreases by $dT_c$, satisfying the relation $dQ_c/dQ_h = T_c/T_h$.
If at $t = 0$ we have $T_c = T_0$, determine the time $t$ that the Carnot
refrigerator must operate for the cooled body to reach a temperature
$T < T_0$. Express your answer in terms of $C_c$, $T_h, P, T_0$ and $T$.

\textbf{T3-C5 [1.5 pts].} The efficiency of a refrigerator is quantified by its
\textit{coefficient of performance} $\mathrm{COP} = Q_c/W$. Determine the
overall COP of the refrigerator of question T3-C4 with all the cycles performed
until time $t$. Write your answer in terms of $T_0$ and $T_h$ and $T$.
